\subsection{Actual prescribed-base native affine atlas}
\label{subsec:relative-native-atlas}

All coefficient actions below are derived from the genuine structure morphism
by the actual section maps. The construction concerns point spaces and their
topological gluing. It supplies no analytic structure sheaf, model-center map,
coherent-sheaf gluing or arbitrary-scheme Hausdorff conclusion.

\begin{theorem}[Genuine section values in the native affine comparison]
\label{scheme:affine-chart-section-values}
\lean{EnvelopeProperty.Berkovich.SchemeValuedPoint.eval_map_ι_topIso,
EnvelopeProperty.Berkovich.SchemeValuedPoint.affineOpenSpectrumHomeomorph_val_apply,
EnvelopeProperty.Berkovich.SchemeValuedPoint.eval_sectionBaseMap}
\leanok
\uses{scheme:native-affine-atlas,scheme:valued-open-embedding,
scheme:affine-valued-point-homeomorphism,scheme:relative-section-coefficients}
The actual native affine-open homeomorphism to the spectrum of $\Gamma(X,U)$
preserves evaluation of every genuine section on $U$. For any native point
supported in $U$, evaluating a derived coefficient from $R$ equals evaluation
of that coefficient at its actual native image under the structure morphism.
\end{theorem}
\begin{proof}
For the open inclusion, the actual inverse section isomorphism and the global
section isomorphism of the open subscheme identify evaluation of $s\in\Gamma(X,U)$
with evaluation of its genuine global section on the open subscheme. The real
scheme affine isomorphism has global section pullback equal to the Spec global
section isomorphism followed by this open global section isomorphism. Their
actual inverse identities therefore identify the affine spectrum coordinate
with the original native value of $s$.

For a derived coefficient, genuine restriction preserves residue evaluation,
and global evaluation naturality through the structure morphism identifies
that value with the base native evaluation. This statement is proved on all
native points, before imposing any coefficient norm condition.
\end{proof}

\begin{theorem}[Actual relative open restrictions]
\label{scheme:relative-open-embedding}
\lean{EnvelopeProperty.Berkovich.RelativeSchemeValuedPoint.supportRelativeHomeomorph,
EnvelopeProperty.Berkovich.RelativeSchemeValuedPoint.openImmersionHomeomorph,
EnvelopeProperty.Berkovich.RelativeSchemeValuedPoint.openImmersionHomeomorph_apply_coe,
EnvelopeProperty.Berkovich.RelativeSchemeValuedPoint.range_map_of_isOpenImmersion,
EnvelopeProperty.Berkovich.RelativeSchemeValuedPoint.isOpenEmbedding_map_of_isOpenImmersion,
EnvelopeProperty.Berkovich.RelativeSchemeValuedPoint.restrictionHomeomorph,
EnvelopeProperty.Berkovich.RelativeSchemeValuedPoint.restrictionHomeomorph_apply_coe}
\leanok
\uses{scheme:valued-open-embedding,scheme:relative-valued-point-functor}
A genuine open immersion commuting with the two structure maps induces an open
embedding of native prescribed-base point spaces. Its image consists exactly
of the fixed-base native target points whose scheme support belongs to the
scheme open image. In particular actual open subschemes give the corresponding
support-open restrictions of the fixed-base native point space.
\end{theorem}
\begin{proof}
The already proved native open-immersion homeomorphism lands on the exact
support-open image. Native functoriality and the actual commuting base diagram
show that its source point has the fixed base point if and only if its target
point does. Restrict this genuine homeomorphism to these exact predicates and
reorder the support and coefficient conditions. The result is the asserted
relative homeomorphism onto the support-open image. Its subtype inclusion is
an open embedding, so the relative point map itself is an open embedding.
No coefficient-field extension or valuation extension is needed.
\end{proof}

\begin{theorem}[Derived-coefficient affine relative chart]
\label{scheme:relative-affine-open-homeomorphism}
\lean{EnvelopeProperty.Berkovich.SchemeValuedPoint.mem_relativeLocus_iff_sectionCoefficients,
EnvelopeProperty.Berkovich.RelativeSchemeValuedPoint.affineOpenSpectrumHomeomorph,
EnvelopeProperty.Berkovich.RelativeSchemeValuedPoint.affineOpenSpectrumHomeomorph_val_apply,
EnvelopeProperty.Berkovich.RelativeSchemeValuedPoint.affineOpenSpectrumHomeomorph_eq_sectionSpectrum,
EnvelopeProperty.Berkovich.RelativeSchemeValuedPoint.affineOpenSpectrumHomeomorph_symm_coe}
\leanok
\uses{scheme:affine-chart-section-values,scheme:relative-open-embedding,
scheme:relative-valued-points,spectrum:relative}
For every genuine affine open $U$ of $X$ with structure map to
$\operatorname{Spec}R$, its fixed-base native support-open subspace is
homeomorphic to $\mathcal S_v(\Gamma(X,U))$, using the actual derived
coefficient action on $\Gamma(X,U)$. Every section value is preserved.
\end{theorem}
\begin{proof}
By the first theorem, the absolute affine chart preserves each derived
coefficient value. A native point maps to the exact fixed base point precisely
when all of its mapped original-ring values agree with $v$, because the
proved affine base homeomorphism is injective and whole seminorm equality is
coordinate equality. This is therefore exactly the prescribed-base predicate
on the affine section-ring spectrum. Restrict the genuine absolute chart
homeomorphism to this pair of predicates, then reorder the support and fixed
base conditions. This constructs both inverse maps and continuity without an
assumed unrelated algebra action or an assumed norm-preservation hypothesis.
\end{proof}

\begin{theorem}[Genuine relative affine point-space gluing]
\label{scheme:relative-native-affine-atlas}
\lean{EnvelopeProperty.Berkovich.RelativeSchemeValuedPoint.supportOpens,
EnvelopeProperty.Berkovich.RelativeSchemeValuedPoint.openCoverGlueData,
EnvelopeProperty.Berkovich.RelativeSchemeValuedPoint.OpenCoverGlued,
EnvelopeProperty.Berkovich.RelativeSchemeValuedPoint.iUnion_supportOpens,
EnvelopeProperty.Berkovich.RelativeSchemeValuedPoint.openCoverGlueHomeomorph,
EnvelopeProperty.Berkovich.RelativeSchemeValuedPoint.affineChartInclusion,
EnvelopeProperty.Berkovich.RelativeSchemeValuedPoint.continuous_affineChartInclusion,
EnvelopeProperty.Berkovich.RelativeSchemeValuedPoint.openCoverGlueHomeomorph_affineChartInclusion,
EnvelopeProperty.Berkovich.RelativeSchemeValuedPoint.isOpenEmbedding_affineChartInclusion,
EnvelopeProperty.Berkovich.RelativeSchemeValuedPoint.affineChartInclusion_eq_iff,
EnvelopeProperty.Berkovich.RelativeSchemeValuedPoint.affineChartInclusion_jointly_surjective,
EnvelopeProperty.Berkovich.RelativeSchemeValuedPoint.canonicalAffineGlueHomeomorph}
\leanok
\uses{scheme:relative-affine-open-homeomorphism,scheme:native-open-cover-gluing,
scheme:canonical-native-affine-atlas}
Any actual scheme affine open cover yields genuine TopCat gluing of the
prescribed-base native support opens, homeomorphic to the fixed-base native
point space. Their pieces are parametrized by the relative spectra of the
actual section rings with their derived coefficient actions. The inclusions
are open embeddings and jointly surjective, and two representatives are
identified exactly when their whole relative native points agree. All actual
affine opens give such a canonical atlas for every scheme over the base.
\end{theorem}
\begin{proof}
Inverse support opens remain open in the fixed-base subspace and still cover
it, since every underlying scheme point lies in a member of the given cover.
Apply the actual topological open-cover gluing construction. Parametrize each
piece by the preceding genuine relative affine homeomorphism. Actual gluing
inclusions are open embeddings; composing with chart inverse homeomorphisms
proves the same for these relative-spectrum parameters. Joint surjectivity
comes from the real gluing inclusions and chart bijections. The actual gluing
homeomorphism maps each representative to its whole relative native point;
injectivity proves the equality criterion. The already proved genuine cover
by all affine opens supplies the canonical case.
\end{proof}
